\documentclass[10pt,a4paper,ssfamily]{exam}
\usepackage[brazil]{babel}
\usepackage[numbers,sort&compress]{natbib}
\usepackage[utf8]{inputenc}
\renewcommand{\baselinestretch}{1.15}
\usepackage{epsfig}
\usepackage{mhchem}
\usepackage{graphicx}
\usepackage{makeidx}
\usepackage{multicol}    
\usepackage{multirow}
\usepackage{amssymb}
\usepackage{amsmath}
\DeclareMathOperator{\sen}{sen}
\newcommand{\listfont}{\bf\fontencoding{OT1}\sffamily\large}
\newcommand{\titlesfont}{\fontencoding{OT1}\sffamily\large}
\renewcommand{\baselinestretch}{1.24}
\newcommand{\Mg}{Mg$^{2+}$}
\renewcommand{\t}{\textrm}
\newcommand{\cc}[1]{[\textrm{#1}]}
\usepackage{enumitem}
\newcommand{\bmat}{\left[\begin{array}{cc}}
\newcommand{\emat}{\end{array}\right]}
\newcommand{\nomera}{
  \noindent
  \begin{tabular}{|p{0.7\textwidth}|p{0.3\textwidth}|}
  Nome: & RA: \\ 
  \hline
  \end{tabular}
}

\newcommand{\qini}{\begin{enumerate}[resume]\item}
\newcommand{\qend}{\end{enumerate}}

\newcommand{\oC}{$^\circ$C}
\newcommand{\kcalmol}{~kcal~mol$^{-1}$}
\newcommand{\moll}{~mol~L$^{-1}$}
\newcommand{\molkg}{~mol~kg$^{-1}$}
\newcommand{\gl}{~g~L$^{-1}$}
\newcommand{\e}[1]{$\times 10^{#1}$}

\newcommand{\eps}{\varepsilon}

\newcommand{\Resposta}[1]{\vspace{0.5cm}\noindent{\bf Resposta:}\\ #1
\begin{center}\rule{\textwidth}{0.4pt}\end{center}}
\renewcommand{\Resposta}[1]{}
\newcommand{\resposta}[1]{\vspace{0.5cm}\noindent{\bf Resposta:}\\ #1
\begin{center}\rule{\textwidth}{0.4pt}\end{center}}
                                                           

\begin{document}
\sffamily
\pagestyle{empty}
\nomera

\begin{center}

{\bf QG101 - Química Geral}

{\bf Aula 1 - Unidades de Medida e Mol}\\
\underline{Justifique suas respostas.}
\end{center}

\begin{center}{\bf Questões}\end{center}

\qini
Na solução de menor energia do rotor rígido quântico, as partículas estão paradas. 
Como o rotor não tem energia potencial, isto significa que a sua energia total é nula.
\begin{enumerate}
\item
Qual a fórmula exata desta solução, com a normalização correta? 
\underline{Mostre} exatamente como obter esta solução.
\item
Faça um gráfico da densidade de probabilidade $Y^2(\theta,\phi)$ em função 
de $\theta$. 
\item
Faça o gráfico qualitativo da probabilidade de encontrar o rotor com um ângulo $\theta\pm\Delta\theta$, em função de $\theta$.
\item
Explique brevemente a relação da densidade de probabilidade ser constante com o princípio de incerteza.
\end{enumerate}
{\bf Obs:} Lembre-se que o arco de círculo que faz ângulo $\theta$ com o eixo $z$
tem comprimento proporcional a $\sen\theta$.
\qend

\qini
As energias das soluções do rotor rígido quântico são dadas pela equação 
\[
E_J = \frac{\hbar^2}{2I}J(J+1)\textrm{ com }J = 0, 1, 2,...
\]
\begin{enumerate}
\item
Desenhe um diagrama de níveis de energia, colocando as energias como função da energia do nível de mais baixa energia.
\item
Desenhe um espectro de absorção indicando as 4 transições de menor energia, assumindo que as transições só podem ocorrer entre níveis de energia consecutivos.
\item
Desenhe um segundo espectro correspondente a um rotor rígido com partículas de maior massa. Deixe claras as diferenças qualitativas entre este espectro e o espectro do item anterior.
\end{enumerate}
\qend

\end{document}

